\documentclass[11pt]{article}
\usepackage[margin=1in]{geometry}
\usepackage{amsmath,amssymb}
\usepackage{enumitem}
\title{COMP3242 Week 2 Practice Exam (Solutions)}
\date{}
\begin{document}
\maketitle

\section*{Q1}
\begin{enumerate}[label=(\alph*)]
  \item $z=a^Tx+b=0.5\cdot1+(-1)\cdot2+0.2=-1.3$.
  \item $\sigma(-1.3)=\frac{1}{1+e^{1.3}}\approx 0.214$.
  \item $L=-\log(0.214)\approx 1.54$.
\end{enumerate}

\section*{Q2}
A single linear classifier defines one hyperplane, so it can only separate data that are linearly separable. XOR is not linearly separable in input space. A two-layer perceptron can create intermediate nonlinear features in the hidden layer, then separate those transformed features in the output layer.

\section*{Q3}
\[
\#A=8\times5=40,\; \#b=8,\; \#C=3\times8=24,\; \#d=3.
\]
Total parameters:
\[
40+8+24+3=75.
\]

\section*{Q4}
\begin{enumerate}[label=(\alph*)]
  \item Softmax + CE: single-label multi-class classification (e.g., classifying an image into exactly one of 10 classes).
  \item Sigmoid + BCE: multi-label classification (e.g., one image may contain multiple attributes simultaneously).
\end{enumerate}
Difference: softmax models mutually exclusive classes (probabilities sum to 1), while sigmoid treats each class as an independent binary decision.

\section*{Q5}
\[
z = xA^T + b: (64,20)(20,50) + (50,) \Rightarrow (64,50).
\]
\[
\text{logits} = zC^T + d: (64,50)(50,10) + (10,) \Rightarrow (64,10).
\]
For `CrossEntropyLoss`, label shape should be `(64,)` with integer class indices in `[0,9]` and dtype `torch.long`.

\end{document}
